% GENERATED FILE. Edit canonical lessons, then run npm run generate:books.
% canonical-source-hash: fnv1a64:115cc633bcbd2a9e
% canonical-lessons: TA-C338-atimai, TA-C338-vitutalai, TA-C338-puratci, TA-C338-kalakam, TA-C338-ciraiyilatai

\chapter{Freedom}
\label{ch:ta-freedom}
\hlchaptermodality{338}

\begin{chapteropening}
\textbf{By the end of this chapter:} \emph{I can talk about struggles for freedom.}

Complete the last lesson of chapter 338: I can talk about struggles for freedom.
\end{chapteropening}

\section[\texorpdfstring{\ta{அடிமை} (aṭimai)}{அடிமை (aṭimai)}]{\ta{அடிமை} (aṭimai) — a slave; slavery}

\label{lesson:TA-C338-atimai}



\begin{quote}
Before the new one: say the Tamil for to deceive, then the Tamil for a disguise.
\end{quote}

\subsection*{You'll want to know: \ta{அடிமை}}
\textbf{\ta{அடிமை}} — \emph{aṭimai} — \textquotedblleft{}a slave; slavery\textquotedblright{}.

\begin{quote}
\ta{பழங்காலத்தில்} \ta{பல} \ta{பேர்} \ta{அடிமைகளாக} \ta{வாழ்ந்தார்கள்}.
\end{quote}

\begin{quote}
\emph{paḻaṅkālattil pala pēr aṭimaikaḷāka vāḻntārkaḷ.}
\end{quote}

\begin{quote}
In ancient times many people lived as slaves.
\end{quote}

Also of a habit: \ta{பழக்கத்துக்கு} \ta{அடிமை}, a slave to a habit.

\begin{cousinweb}
\textbf{\ta{அடிமை}} (\emph{aṭimai}) is native Dravidian: DEDR 72 gives it under \ta{அடி}, \textquotedblleft{}foot\textquotedblright{}, one who is at the feet of another.
\end{cousinweb}

\begin{grammarlens}[title={What you've built}]
One more word for telling what happened.
\end{grammarlens}

\subsection*{Your turn}
\begin{itemize}
\raggedright
  \item \emph{Say it:} \emph{aṭimai}
  \item \emph{Say it:} \emph{aṭimai}, once more
  \item \emph{Say it:} \emph{paḻaṅkālattil pala pēr aṭimaikaḷāka vāḻntārkaḷ.} — In ancient times many people lived as slaves.
  \item \emph{Recall:} say the Tamil for to deceive, then the Tamil for a disguise, then say \emph{aṭimai} again
\end{itemize}

\begin{tcolorbox}[breakable,colback=teal!4,colframe=teal!35!black,title={Before you move on}]
What does \ta{அடிமை} mean? (\textquotedblleft{}A slave; slavery\textquotedblright{}.) Say it once more.
\end{tcolorbox}
\section[\texorpdfstring{\ta{விடுதலை} (viṭutalai)}{விடுதலை (viṭutalai)}]{\ta{விடுதலை} (viṭutalai) — liberation; freedom; release}

\label{lesson:TA-C338-vitutalai}



\begin{quote}
Before the new one: say the Tamil for a disguise, then the Tamil for a slave.
\end{quote}

\subsection*{You'll want to know: \ta{விடுதலை}}
\textbf{\ta{விடுதலை}} — \emph{viṭutalai} — \textquotedblleft{}liberation; freedom; release\textquotedblright{}.

\begin{quote}
\ta{நீண்ட} \ta{போராட்டத்துக்குப்} \ta{பிறகு} \ta{நாட்டுக்கு} \ta{விடுதலை} \ta{கிடைத்தது}.
\end{quote}

\begin{quote}
\emph{nīṇṭa pōrāṭṭattukkup piṟaku nāṭṭukku viṭutalai kiṭaittatu.}
\end{quote}

\begin{quote}
After a long struggle the country gained its freedom.
\end{quote}

Also release from prison: \ta{கைதிக்கு} \ta{விடுதலை}.

\begin{cousinweb}
\textbf{\ta{விடுதலை}} (\emph{viṭutalai}) is native Dravidian: a noun from \ta{விடு}, \textquotedblleft{}to leave, let go\textquotedblright{} (DEDR 5393).
\end{cousinweb}

\begin{grammarlens}[title={What you've built}]
One more word for telling what happened.
\end{grammarlens}

\subsection*{Your turn}
\begin{itemize}
\raggedright
  \item \emph{Say it:} \emph{viṭutalai}
  \item \emph{Say it:} \emph{viṭutalai}, once more
  \item \emph{Say it:} \emph{nīṇṭa pōrāṭṭattukkup piṟaku nāṭṭukku viṭutalai kiṭaittatu.} — After a long struggle the country gained its freedom.
  \item \emph{Recall:} say the Tamil for a disguise, then the Tamil for a slave, then say \emph{viṭutalai} again
\end{itemize}

\begin{tcolorbox}[breakable,colback=teal!4,colframe=teal!35!black,title={Before you move on}]
What does \ta{விடுதலை} mean? (\textquotedblleft{}Liberation; freedom; release\textquotedblright{}.) Say it once more.
\end{tcolorbox}
\section[\texorpdfstring{\ta{புரட்சி} (puraṭci)}{புரட்சி (puraṭci)}]{\ta{புரட்சி} (puraṭci) — a revolution}

\label{lesson:TA-C338-puratci}



\begin{quote}
Before the new one: say the Tamil for a slave, then the Tamil for liberation.
\end{quote}

\subsection*{You'll want to know: \ta{புரட்சி}}
\textbf{\ta{புரட்சி}} — \emph{puraṭci} — \textquotedblleft{}a revolution\textquotedblright{}.

\begin{quote}
\ta{அந்தப்} \ta{புரட்சி} \ta{நாட்டின்} \ta{வரலாற்றை} \ta{மாற்றியது}.
\end{quote}

\begin{quote}
\emph{antap puraṭci nāṭṭiṉ varalāṟṟai māṟṟiyatu.}
\end{quote}

\begin{quote}
That revolution changed the country's history.
\end{quote}

The adjective is \ta{புரட்சிகரமான}, revolutionary.

\begin{cousinweb}
\textbf{\ta{புரட்சி}} (\emph{puraṭci}) is native Dravidian: DEDR 4285 gives it under \ta{புரள்}, \textquotedblleft{}to be upset, roll over\textquotedblright{}, as \textquotedblleft{}upsetting, disorder\textquotedblright{}.
\end{cousinweb}

\begin{grammarlens}[title={What you've built}]
One more word for telling what happened.
\end{grammarlens}

\subsection*{Your turn}
\begin{itemize}
\raggedright
  \item \emph{Say it:} \emph{puraṭci}
  \item \emph{Say it:} \emph{puraṭci}, once more
  \item \emph{Say it:} \emph{antap puraṭci nāṭṭiṉ varalāṟṟai māṟṟiyatu.} — That revolution changed the country's history.
  \item \emph{Recall:} say the Tamil for a slave, then the Tamil for liberation, then say \emph{puraṭci} again
\end{itemize}

\begin{tcolorbox}[breakable,colback=teal!4,colframe=teal!35!black,title={Before you move on}]
What does \ta{புரட்சி} mean? (\textquotedblleft{}A revolution\textquotedblright{}.) Say it once more.
\end{tcolorbox}
\section[\texorpdfstring{\ta{கலகம்} (kalakam)}{கலகம் (kalakam)}]{\ta{கலகம்} (kalakam) — a revolt; a quarrel}

\label{lesson:TA-C338-kalakam}



\begin{quote}
Before the new one: say the Tamil for liberation, then the Tamil for a revolution.
\end{quote}

\subsection*{You'll want to know: \ta{கலகம்}}
\textbf{\ta{கலகம்}} — \emph{kalakam} — \textquotedblleft{}a revolt; a quarrel\textquotedblright{}.

\begin{quote}
\ta{வரியை} \ta{எதிர்த்து} \ta{மக்கள்} \ta{கலகம்} \ta{செய்தார்கள்}.
\end{quote}

\begin{quote}
\emph{variyai etirttu makkaḷ kalakam ceytārkaḷ.}
\end{quote}

\begin{quote}
The people revolted against the tax.
\end{quote}

\ta{கலகம்} \ta{செய்}, to revolt.

\begin{cousinweb}
\textbf{\ta{கலகம்}} (\emph{kalakam}) is borrowed from Sanskrit \emph{kalaha}, \textquotedblleft{}quarrel, fight\textquotedblright{} (CDIAL 2922).
\end{cousinweb}

\begin{grammarlens}[title={What you've built}]
One more word for telling what happened.
\end{grammarlens}

\subsection*{Your turn}
\begin{itemize}
\raggedright
  \item \emph{Say it:} \emph{kalakam}
  \item \emph{Say it:} \emph{kalakam}, once more
  \item \emph{Say it:} \emph{variyai etirttu makkaḷ kalakam ceytārkaḷ.} — The people revolted against the tax.
  \item \emph{Recall:} say the Tamil for liberation, then the Tamil for a revolution, then say \emph{kalakam} again
\end{itemize}

\begin{tcolorbox}[breakable,colback=teal!4,colframe=teal!35!black,title={Before you move on}]
What does \ta{கலகம்} mean? (\textquotedblleft{}A revolt; a quarrel\textquotedblright{}.) Say it once more.
\end{tcolorbox}
\section[\texorpdfstring{\ta{சிறையிலடை} (ciṟaiyilaṭai)}{சிறையிலடை (ciṟaiyilaṭai)}]{\ta{சிறையிலடை} (ciṟaiyilaṭai) — to jail; to imprison}

\label{lesson:TA-C338-ciraiyilatai}



\begin{quote}
Before the new one: say the Tamil for a revolution, then the Tamil for a revolt.
\end{quote}

\subsection*{You'll want to know: \ta{சிறையிலடை}}
\textbf{\ta{சிறையிலடை}} — \emph{ciṟaiyilaṭai} — \textquotedblleft{}to jail; to imprison\textquotedblright{}.

\begin{quote}
\ta{அரசு} \ta{அந்தத்} \ta{தலைவர்களைச்} \ta{சிறையிலடைத்தது}.
\end{quote}

\begin{quote}
\emph{aracu antat talaivarkaḷaic ciṟaiyilaṭaittatu.}
\end{quote}

\begin{quote}
The government jailed those leaders.
\end{quote}

Past (I …): \ta{சிறையிலடைத்தேன்}, \emph{ciṟaiyilaṭaittēṉ}.

\ta{சிறையில்}, in prison + \ta{அடை}, to shut in.

\begin{cousinweb}
\textbf{\ta{சிறையிலடை}} (\emph{ciṟaiyilaṭai}) is a Tamil compound built from \ta{சிறையில்}, \textquotedblleft{}in prison\textquotedblright{}, from \ta{சிறை}, \textquotedblleft{}confinement\textquotedblright{} (DEDR 1980), and \ta{அடை}, \textquotedblleft{}to shut, lock\textquotedblright{} (DEDR 83).
\end{cousinweb}

\begin{grammarlens}[title={What you've built}]
One more word for telling what happened.
\end{grammarlens}

\subsection*{Your turn}
\begin{itemize}
\raggedright
  \item \emph{Say it:} \emph{ciṟaiyilaṭai}
  \item \emph{Say it:} \emph{ciṟaiyilaṭai}, once more
  \item \emph{Say it:} \emph{aracu antat talaivarkaḷaic ciṟaiyilaṭaittatu.} — The government jailed those leaders.
  \item \emph{Recall:} say the Tamil for a revolution, then the Tamil for a revolt, then say \emph{ciṟaiyilaṭai} again
\end{itemize}

\begin{tcolorbox}[breakable,colback=teal!4,colframe=teal!35!black,title={Before you move on}]
What does \ta{சிறையிலடை} mean? (\textquotedblleft{}To jail; to imprison\textquotedblright{}.) Say it once more.
\end{tcolorbox}
